\documentclass[11pt]{article}

\input{macros}

\title{SHARP Reproduction: Methods and Local-Compute Algorithmic Plan}
\author{}
\date{\today}

\begin{document}
\maketitle

\begin{abstract}
This note summarizes the SHARP reproduction target \cite{gross2025sharp} and
the algorithmic changes needed to make the computation practical on a local
workstation. The goal is not to change the scientific claim during
reproduction, but to reproduce the published network-proximity and pAGE
computations, then replace the slow parts with provably equivalent
implementations. The main point is that rented cloud compute should not be
necessary: the expensive step is mostly repeated drug-independent work inside a
permutation null, and that work can be shared, vectorized, and checked against a
naive oracle.
\end{abstract}

\section{Objects and Data}

Let \(G=(V,E)\) be the undirected human protein interaction graph, with shortest
path distance \(d(u,v)\). Vertices are genes or proteins after identifier
harmonization. For each aging hallmark \(h\), let \(M_h \subset V\) be its
OpenGenes-derived hallmark module, possibly filtered by confidence threshold
\(\ell \in \{1,\ldots,5\}\). For each drug \(a\), let \(T_a \subset V\) be the
set of its mapped protein targets. Let \(A_g \in \{-1,0,1\}\) denote the
age-associated direction for gene \(g\), with \(+1\) meaning increased with age
and \(-1\) meaning decreased with age. Let \(D_{a,g} \in \{-1,0,1\}\) denote
the drug-induced expression direction from CMap/LINCS.

The local reproduction inputs are the processed interactome, hallmark
memberships, age-direction table, DrugBank-derived target mappings, CMap
metadata, published proximity/pAGE result tables, and the Level 5 CMap GCTX
matrix. Before recomputing pAGE from the GCTX, require a file-size gate:
\[
  \operatorname{size}(\texttt{level5\_beta\_trt\_cp\_n720216x12328.gctx})
  = 35{,}518{,}405{,}386 \text{ bytes}.
\]
The current local file must be treated as incomplete until this equality holds.

\section{Baseline SHARP Method}

\paragraph{Hallmark modules.}
Each hallmark is represented by the induced subgraph on \(M_h\). The
connectivity check asks whether the largest connected component of \(M_h\) is
larger than expected under degree-matched random gene sets of the same size.
This establishes that each hallmark is not merely a list of genes, but a
network-localized module. The reproduction must preserve the authors'
confidence-level and de-duplication rules, because headline counts depend on
them.

\paragraph{Network proximity.}
SHARP uses a Guney-style closest-distance score between drug targets \(T_a\)
and a hallmark module \(M_h\):
\[
  c_{M_h}(t) = \min_{m \in M_h} d(t,m),
  \qquad
  \delta(a,h) = \frac{1}{|T_a|}\sum_{t \in T_a} c_{M_h}(t).
\]
Smaller \(\delta(a,h)\) means the drug targets lie closer to the hallmark
module. Statistical significance is computed from a degree-matched permutation
null. If \(\delta^{(1)},\ldots,\delta^{(R)}\) are null proximities, then
\[
  z(a,h)=
  \frac{\delta(a,h)-\widehat{\mu}_{a,h}}{\widehat{\sigma}_{a,h}},
  \qquad
  \widehat{\mu}_{a,h}=\frac{1}{R}\sum_{r=1}^R \delta^{(r)}, \quad
  \widehat{\sigma}_{a,h}^2=\frac{1}{R-1}\sum_{r=1}^R
  \bigl(\delta^{(r)}-\widehat{\mu}_{a,h}\bigr)^2.
\]
The published significance rule is \(z<-1.96\), with marginal calls sometimes
reported separately. The key unresolved reproduction detail is whether the null
randomizes only the drug target set, or both \(T_a\) and \(M_h\). This must be
verified from the author code before full-scale execution.

\paragraph{Expression directionality: pAGE.}
pAGE is a signed concordance score between age-associated expression and
drug-induced expression:
\[
  \operatorname{pAGE}(a;H)
  = -\sum_{g \in H} A_g D_{a,g}.
\]
Here \(H\) is the chosen gene universe, typically the genes with usable age and
CMap directions. The leading minus sign makes reversal positive: if a gene
increases with age but the drug decreases it, then \(A_gD_{a,g}=-1\), which
contributes \(+1\). Thus \(\operatorname{pAGE}>0\) is pro-longevity by the
paper's convention, while \(\operatorname{pAGE}<0\) is age-accelerating.

\paragraph{Combined candidate rule.}
The baseline candidate screen is conceptually two-stage. First, a drug must be
network-proximal to at least one hallmark module under the proximity null.
Second, among drugs with CMap coverage, pAGE must be positive. The reproduction
should report counts at each filtering step: genes, modules, drugs, proximal
drug--hallmark pairs, CMap-covered proximal drugs, and pAGE-positive candidates.

\section{Why the Naive Computation Is Slow}

The observed proximity scores are cheap. The expensive part is the permutation
null: for each drug, hallmark, confidence level, and randomization, the naive
implementation repeatedly evaluates
\(\min_{m\in M_h} d(t,m)\) in Python-level loops. If \(N_d\) drugs,
\(N_h\) hallmark/level combinations, and \(R\approx 10^3\) permutations are
used, this produces on the order of \(N_dN_hR\) proximity evaluations. The
problem is not that shortest paths are intrinsically too expensive; it is that
the same module-distance information is recomputed or re-read in the wrong
place.

\section{Recommended Local-Compute Algorithm}

\paragraph{Correctness anchor.}
First implement the naive permutation method on a small deterministic subset.
This oracle should use the same random seed, degree bins, confidence filters,
and missing-gene rules as the authors. All optimized implementations must agree
with the oracle to numerical tolerance, for example
\[
  |z_{\mathrm{fast}}(a,h)-z_{\mathrm{oracle}}(a,h)| < 10^{-9}
\]
on the held-out subset. Reproduction fidelity comes before speed.

\paragraph{Multi-source BFS.}
For a fixed module \(M_h\), compute the whole vector
\[
  c_{M_h} = \bigl(c_{M_h}(v)\bigr)_{v\in V}
\]
with one multi-source breadth-first search seeded at all nodes in \(M_h\). Then
each drug proximity is only a gather-and-mean:
\[
  \delta(a,h) = \operatorname{mean}\{c_{M_h}(t):t\in T_a\}.
\]
This removes the graph traversal and the module minimum from the inner drug
loop.

\paragraph{Share the null.}
If the null randomizes hallmark modules, precompute \(R\) module-distance
vectors for each hallmark/level/network combination and store them as an
\(R\times |V|\) integer array. Most interactome distances are small, so
\texttt{int8} or \texttt{int16} is usually sufficient after reserving one value
for unreachable nodes. For every drug, the null distribution is obtained by
gathering the target columns from this matrix and taking row means. If the null
randomizes drug targets only, the savings are even larger: the observed module
distance vector is reused, and random target draws become a vectorized indexed
mean over \(c_{M_h}\).

\paragraph{Avoid all-pairs distances.}
Do not load or regenerate a dense all-pairs distance pickle unless an explicit
storage estimate justifies it. A full \(n\times n\) distance matrix is the wrong
object for this task. The needed object is a small family of one-to-all
distance vectors \(c_{M_h}\), plus the corresponding random-module vectors if
the null requires them.

\paragraph{Stream CMap, do not materialize more than needed.}
For pAGE, the Level 5 GCTX is large but not algorithmically hard. Once the file
is complete, read only the signatures and genes needed for the reproduction,
map them through \texttt{siginfo\_beta.txt} and \texttt{geneinfo\_beta.txt},
and compute \(D_{a,g}\) or continuous drug effects in chunks. The expression
step should be I/O-bound and memory-bounded, not a cloud-scale computation.

\section{Methodological Improvements After Reproduction}

\paragraph{Effect-size-aware reversal.}
pAGE discards effect size and uncertainty. After exact reproduction, replace it
with a continuous reversal score
\[
  S(a,h,c) =
  -\sum_{g\in H_h} w_g\,\beta_{\mathrm{age}}(g,c)\,
  \beta_{\mathrm{drug}}(a,g,c),
\]
where \(c\) denotes biological context, such as tissue, cell line, dose, or
time. The sign convention matches pAGE: larger is more age-reversing. This
upgrade is only half available immediately. CMap supplies continuous drug-side
effects, but the current OpenGenes age-side input is categorical, so a fully
continuous \(\beta_{\mathrm{age}}\) requires a quantitative aging expression
resource.

\paragraph{Context sensitivity.}
Do not collapse dose, exposure time, and cell line prematurely. At minimum,
report robustness across CMap contexts. A simple model is
\[
  \beta_{\mathrm{drug}}(a,g,c)
  = \alpha_{a,g}+\gamma_{\mathrm{cell}(c),g}
    + \eta_{\mathrm{dose}(c),g}+\tau_{\mathrm{time}(c),g}
    + \varepsilon_{a,g,c},
\]
or a robust nonparametric aggregation if sample sizes are too thin for a mixed
model. Candidate drugs whose sign changes across contexts should be marked
fragile.

\paragraph{Better validation.}
Validation should include ablations: proximity only, pAGE only, target degree,
target overlap, shuffled age signatures, shuffled drug targets, and
degree-preserving network controls. Because there is no clean negative set for
``does not affect longevity,'' precision estimates must be interpreted
cautiously. Decoy and permutation nulls are safer than pretending unknown drugs
are true negatives.

\section{Feasibility Claim}

The reproduction should not require AWS if three conditions hold:
\begin{enumerate}[leftmargin=*]
  \item the Level 5 GCTX download is complete and verified by byte size or
  checksum;
  \item the proximity implementation is validated against a naive oracle before
  full-scale execution;
  \item the full run uses module-distance vectors and vectorized null
  evaluation rather than Python loops or dense all-pairs distance matrices.
\end{enumerate}
Under these conditions, proximity should move from a weeks-scale script to a
hours-scale local job, and pAGE should be a streaming matrix-processing task.
The remaining hard problems are methodological validation and data
harmonization, not raw compute rental.

\bibliographystyle{alpha}
\bibliography{references}

\end{document}
